%!TEX root=main.tex

\section{Introduction}\label{sect:intro}
\EZComm{DA FARE, intanto ho messo questo schemino}
Due idee nuove: 

1) Uso di risorse indiretto con viewpoint adaptation. Ossia, se una risorsa contiene un riferimento ad un'altra risorsa, quando accedo attraverso la prima consumo la seconda, ma in un modo mediato dal contenitore. 

2) Gradi assunzione

Nel nostro caso Java-like, l'annotazione di uso del campo indica come modifico il consumo del reference attraverso il receiver. In altri termini, per ogni campo mi dice quanto posso usare quel campo "passando da" quel receiver. Formalmente, quando accedo a un campo consumo il reference corrispondente, ma il risultato si ottiene tramite prodotto con il grado del receiver.  Nel caso di gradi non quantitativi, per esempio 0/readonly/mut, questo corrisponde a viewpoint adaptation. 

We illustrate our approach on a tiny imperative Java-like language, in the spirit of the seminal calculus Featherweight Java (FJ) \cite{IgarashiPW99}. 
Besides the FJ constructs, the language includes some imperative ones, notably field update and deallocation, and evaluation takes place in a memory, a map from references (unique object identifiers) into object states. 
\EZComm{ci sono stati diversi calcoli imperativi in letteratura forse citarli}


